% GENERATED FILE. Edit canonical lessons, then run npm run generate:books.
% canonical-source-hash: fnv1a64:9332d0a1c41d8e2a
% canonical-lessons: TA-C106-takecare, TA-C106-safely, TA-C106-late, TA-C106-letsgo, TA-C106-holiday

\chapter{Take Care, Go Safely}
\label{ch:ta-take-care}
\hlchaptermodality{106}

\begin{chapteropening}
\textbf{By the end of this chapter:} \emph{I can tell someone to take care and go safely, say it is late, suggest we go, and talk about a holiday.}

Complete the last lesson of chapter 106: I can tell someone to take care and go safely, say it is late, suggest we go, and talk about a holiday.
\end{chapteropening}

\section[\texorpdfstring{\ta{பார்த்துக்கொள்ளுங்கள்} (pārttukkoḷḷuṅkaḷ)}{பார்த்துக்கொள்ளுங்கள் (pārttukkoḷḷuṅkaḷ)}]{\ta{பார்த்துக்கொள்ளுங்கள்} (pārttukkoḷḷuṅkaḷ) — take care}

\label{lesson:TA-C106-takecare}



\begin{quote}
Before the new one: say the Tamil for a college, then the Tamil for a school.
\end{quote}

\subsection*{You'll want to know: \ta{பார்த்துக்கொள்ளுங்கள்}}
\textbf{\ta{பார்த்துக்கொள்ளுங்கள்}} (\emph{pārttukkoḷḷuṅkaḷ}) — \textquotedblleft{}take care\textquotedblright{}, literally \textquotedblleft{}look after (yourself)\textquotedblright{}.

It is built on \textbf{\ta{பார்}}, \textquotedblleft{}to see, to look\textquotedblright{}.

\textbf{In the family, and next door}

\noindent
\begin{tabularx}{\linewidth}{@{}>{\raggedright\arraybackslash}X>{\raggedright\arraybackslash}X@{}}
\toprule
\textbf{} & \textbf{word} \\
\midrule
Kannada & \ka{ಹುಷಾರು} (\emph{huṣāru}) \\
Malayalam & \ml{സൂക്ഷിക്കണേ} (\emph{sūkṣikkaṇē}) \\
Hindi (neighbour) & \dv{ध्यान} \dv{रखिए} (\emph{dhyān rakhie}) \\
English & take care \\
\bottomrule
\end{tabularx}

\begin{grammarlens}[title={What you've built}]
The kindest goodbye.
\end{grammarlens}

\subsection*{Your turn}
\begin{itemize}
\raggedright
  \item \emph{Say it:} \emph{pārttukkoḷḷuṅkaḷ}
  \item \emph{Say it:} \emph{pārttukkoḷḷuṅkaḷ}, once more
  \item \emph{Say it:} say \emph{pārttukkoḷḷuṅkaḷ}
  \item \emph{Recall:} say the Tamil for a college, then the Tamil for a school, then say \emph{pārttukkoḷḷuṅkaḷ} again
\end{itemize}

\begin{tcolorbox}[breakable,colback=teal!4,colframe=teal!35!black,title={Before you move on}]
What does \ta{பார்த்துக்கொள்ளுங்கள்} mean? (\textquotedblleft{}Take care\textquotedblright{}.) Say it once more.
\end{tcolorbox}
\section[\texorpdfstring{\ta{பத்திரமாக} (pattiramāka)}{பத்திரமாக (pattiramāka)}]{\ta{பத்திரமாக} (pattiramāka) — safely}

\label{lesson:TA-C106-safely}



\begin{quote}
Before the new one: say the Tamil for a school, then the Tamil for take care.
\end{quote}

\subsection*{You'll want to know: \ta{பத்திரமாக}}
\textbf{\ta{பத்திரமாக}} (\emph{pattiramāka}) — \textquotedblleft{}safely\textquotedblright{}.

\textbf{\ta{பத்திரமாகப்} \ta{போங்கள்}} — \textquotedblleft{}go safely\textquotedblright{}.

\begin{grammarlens}[title={What you've built}]
A wish for the road.
\end{grammarlens}

\subsection*{Your turn}
\begin{itemize}
\raggedright
  \item \emph{Say it:} \emph{pattiramāka}
  \item \emph{Say it:} \emph{pattiramāka}, once more
  \item \emph{Say it:} say \emph{pattiramākap pōṅkaḷ}
  \item \emph{Recall:} say the Tamil for a school, then the Tamil for take care, then say \emph{pattiramāka} again
\end{itemize}

\begin{tcolorbox}[breakable,colback=teal!4,colframe=teal!35!black,title={Before you move on}]
What does \ta{பத்திரமாக} mean? (\textquotedblleft{}Safely\textquotedblright{}.) Say it once more.
\end{tcolorbox}
\section[\texorpdfstring{\ta{நேரமாகிவிட்டது} (nēramākiviṭṭatu)}{நேரமாகிவிட்டது (nēramākiviṭṭatu)}]{\ta{நேரமாகிவிட்டது} (nēramākiviṭṭatu) — it has got late}

\label{lesson:TA-C106-late}



\begin{quote}
Before the new one: say the Tamil for take care, then the Tamil for safely.
\end{quote}

\subsection*{You'll want to know: \ta{நேரமாகிவிட்டது}}
\textbf{\ta{நேரமாகிவிட்டது}} (\emph{nēramākiviṭṭatu}) — \textquotedblleft{}it has got late\textquotedblright{}, literally \textquotedblleft{}it has become time\textquotedblright{}.

\textbf{\ta{நேரம்}} is \textquotedblleft{}time\textquotedblright{}. The first reason to leave.

\textbf{In the family, and next door}

\noindent
\begin{tabularx}{\linewidth}{@{}>{\raggedright\arraybackslash}X>{\raggedright\arraybackslash}X@{}}
\toprule
\textbf{} & \textbf{word} \\
\midrule
Hindi (neighbour) & \dv{देर} \dv{हो} \dv{गई} (\emph{der ho gaī}) \\
English & it has got late \\
\bottomrule
\end{tabularx}

\begin{grammarlens}[title={What you've built}]
A reason to go.
\end{grammarlens}

\subsection*{Your turn}
\begin{itemize}
\raggedright
  \item \emph{Say it:} \emph{nēramākiviṭṭatu}
  \item \emph{Say it:} \emph{nēramākiviṭṭatu}, once more
  \item \emph{Say it:} say \emph{nēramākiviṭṭatu}
  \item \emph{Recall:} say the Tamil for take care, then the Tamil for safely, then say \emph{nēramākiviṭṭatu} again
\end{itemize}

\begin{tcolorbox}[breakable,colback=teal!4,colframe=teal!35!black,title={Before you move on}]
What does \ta{நேரமாகிவிட்டது} mean? (\textquotedblleft{}It has got late\textquotedblright{}.) Say it once more.
\end{tcolorbox}
\section[\texorpdfstring{\ta{போகலாம்} (pōkalām)}{போகலாம் (pōkalām)}]{\ta{போகலாம்} (pōkalām) — let's go}

\label{lesson:TA-C106-letsgo}



\begin{quote}
Before the new one: say the Tamil for safely, then the Tamil for it has got late.
\end{quote}

\subsection*{You'll want to know: \ta{போகலாம்}}
\textbf{\ta{போகலாம்}} (\emph{pōkalām}) — \textquotedblleft{}let's go\textquotedblright{}, \textquotedblleft{}we can go\textquotedblright{}.

It is \textbf{\ta{போ}}, \textquotedblleft{}go\textquotedblright{}, with \textbf{-\ta{லாம்}}, \textquotedblleft{}may, let us\textquotedblright{}.

\textbf{In the family, and next door}

\noindent
\begin{tabularx}{\linewidth}{@{}>{\raggedright\arraybackslash}X>{\raggedright\arraybackslash}X>{\raggedright\arraybackslash}X@{}}
\toprule
\textbf{} & \textbf{word} & \textbf{same root} \\
\midrule
Kannada & \ka{ಹೋಗೋಣ} (\emph{hōgōṇa}) &  \\
Telugu & \te{పదండి} (\emph{padaṇḍi}) &  \\
Malayalam & \ml{പോകാം} (\emph{pōkāṁ}) & yes \\
Hindi (neighbour) & \dv{चलिए} (\emph{calie}) &  \\
English & let's go &  \\
\bottomrule
\end{tabularx}

\begin{grammarlens}[title={What you've built}]
Late, so let's go.
\end{grammarlens}

\subsection*{Your turn}
\begin{itemize}
\raggedright
  \item \emph{Say it:} \emph{pōkalām}
  \item \emph{Say it:} \emph{pōkalām}, once more
  \item \emph{Say it:} say \emph{pōkalām}
  \item \emph{Recall:} say the Tamil for safely, then the Tamil for it has got late, then say \emph{pōkalām} again
\end{itemize}

\begin{tcolorbox}[breakable,colback=teal!4,colframe=teal!35!black,title={Before you move on}]
What does \ta{போகலாம்} mean? (\textquotedblleft{}Let's go\textquotedblright{}.) Say it once more.
\end{tcolorbox}
\section[\texorpdfstring{\ta{விடுமுறை} (viṭumuṟai)}{விடுமுறை (viṭumuṟai)}]{\ta{விடுமுறை} (viṭumuṟai) — a holiday}

\label{lesson:TA-C106-holiday}



\begin{quote}
Before the new one: say the Tamil for it has got late, then the Tamil for let's go.
\end{quote}

\subsection*{You'll want to know: \ta{விடுமுறை}}
\textbf{\ta{விடுமுறை}} (\emph{viṭumuṟai}) — \textquotedblleft{}a holiday\textquotedblright{}, \textquotedblleft{}leave\textquotedblright{}.

\textbf{\ta{இன்று} \ta{விடுமுறை}} — \textquotedblleft{}today is a holiday\textquotedblright{}.

\textbf{In the family, and next door}

\noindent
\begin{tabularx}{\linewidth}{@{}>{\raggedright\arraybackslash}X>{\raggedright\arraybackslash}X@{}}
\toprule
\textbf{} & \textbf{word} \\
\midrule
Kannada & \ka{ರಜೆ} (\emph{raje}) \\
Malayalam & \ml{അവധി} (\emph{avadhi}) \\
Hindi (neighbour) & \dv{छुट्टी} (\emph{chuṭṭī}) \\
English & a holiday \\
\bottomrule
\end{tabularx}

\begin{grammarlens}[title={What you've built}]
Take care, safely, it's late, let's go, a holiday.
\end{grammarlens}

\subsection*{Your turn}
\begin{itemize}
\raggedright
  \item \emph{Say it:} \emph{viṭumuṟai}
  \item \emph{Say it:} \emph{viṭumuṟai}, once more
  \item \emph{Say it:} say \emph{iṉṟu viṭumuṟai}
  \item \emph{Recall:} say the Tamil for it has got late, then the Tamil for let's go, then say \emph{viṭumuṟai} again
\end{itemize}

\begin{tcolorbox}[breakable,colback=teal!4,colframe=teal!35!black,title={Before you move on}]
What does \ta{விடுமுறை} mean? (\textquotedblleft{}A holiday\textquotedblright{}.) Say it once more.
\end{tcolorbox}
